\documentclass[11pt]{article}
\usepackage{microtype,amsmath,amssymb,amsthm,mathtools,booktabs,array}
\usepackage[margin=1.05in]{geometry}
\usepackage[colorlinks=true,linkcolor=blue!45!black,citecolor=blue!45!black,urlcolor=blue!45!black]{hyperref}
\usepackage{xcolor}
\hypersetup{pdftitle={The dimension-degree boundary for Gaussian moment vanishing},pdfsubject={Working mathematical manuscript}}
\newtheorem{theorem}{Theorem}[section]
\newtheorem{proposition}[theorem]{Proposition}
\newtheorem{lemma}[theorem]{Lemma}
\newtheorem{corollary}[theorem]{Corollary}
\theoremstyle{definition}\newtheorem{definition}[theorem]{Definition}
\theoremstyle{remark}\newtheorem{remark}[theorem]{Remark}
\newcommand{\E}{\mathbb E}
\newcommand{\C}{\mathbb C}
\newcommand{\tr}{\operatorname{tr}}
\newcommand{\wt}{\operatorname{wt}}
\newcommand{\G}{\mathsf G}
\DeclareMathOperator{\Span}{span}
\title{Punctured polynomial phases and logarithmic Gaussian transport}
\author{}\date{17 November 2026: retained written proof}
\begin{document}\maketitle
This is a written proof over a valued field, including arbitrary
nonconvergent formal Laurent coefficients. It is not an all-ring
statement. It gives an arbitrary-degree logarithmic-volume obstruction and
its formal Gaussian transport. The polynomial-frame lemmas needed
for the proof are written out below. Priority is unassessed;
independent replay, machine-formal verification and external review
are not claimed. No finite computation proves its completeness lemmas.

\section{Gaussian conventions for the applications}
Let $T$ be standard Gaussian and $Z,W$ circular coordinates with
$\E[Z^iW^j]=\delta_{ij}i!$, independent of $T$. For a polynomial
$P$, set $\Phi_Q=\sum_{m\geq0}\E[QP^m]t^m/m!$. Purity means
$\Phi_1=1$; safety means every $\Phi_Q$ is a polynomial.
These are formal series, without an integral-convergence hypothesis.
Gaussian integration by parts is $\E[(\partial_T-T)H]=0$.

For any polynomial $h(z,T)$, let $g=th(g,T)$ be the unique small
formal solution. The exact circular identity is
\[
 \E_{Z,W}[A(Z,T)e^{tWh(Z,T)}]
       =\frac{A(g,T)}{1-th_Z(g,T)}.
\]
Indeed the left side is $\sum_m t^m[z^m]A(z,T)h(z,T)^m$.
It is the formal residue of $A(z,T)/(z-th(z,T))$, whose unique
small pole is $g$. This also proves the identity for a
$t$-dependent polynomial-coefficient formal numerator $A$.

\section{The full formal frame and its Gaussian applications}
\subsection{A formal leading-coefficient criterion}\label{sec:stationaryballs}

For a polynomial phase, the remaining shrinking argument can be
replaced by a valuation calculation. The construction below includes
critical points that collide, remain multiple, or escape to infinity.
The key is to group them before making a Gaussian approximation.

Write $v$ for the valuation of a formal parameter $r$, normalized by
$v(r)=1$. The valuation of a polynomial is the least valuation of its
coefficients. We use the formal Newton--Puiseux theorem: finitely many
algebraic elements over $\C((r))$ belong to one field
$\C((r^{1/N}))$. Convergence is not required~\cite[Chapter~2]{Wall}.

\begin{lemma}[Polynomial separation of formal balls]\label{lem:ballseparation}
Let finitely many disjoint balls have centers $c_B$ in a Puiseux field
and radii specified by $v(x-c_B)\geq\rho_B$. After one common
ramification, put $\delta_B=r^{\rho_B}$. For any one ball $B$ and
integer $j\geq0$, there is a polynomial $q_{B,j}(x)$ over the Laurent
field whose substitution $x=c_C+\delta_C y$ has integral coefficients
for every ball $C$, and whose reduction is $y^j$ on $B$ and zero on
every other ball.
\end{lemma}
\begin{proof}
The finitely many valuations $v(c_B-c_C)$ give a finite tree of
clusters of centers. At a node with scale $r^d$, the coordinate
$u=(x-c)/r^d$ has distinct constant reductions $a_1,\ldots,a_s$
on its children. The Lagrange polynomial
\[
 e_i(u)=\prod_{k\ne i}\frac{u-a_k}{a_i-a_k}
\]
reduces to one on child $i$ and zero on every other child.

Start at the chosen leaf with $((x-c_B)/\delta_B)^j$. Induct upward.
Suppose a polynomial $h$ already has the required reductions on the
chosen child. On each of the other child balls its coefficients have
some finite lower valuation. The restriction of $e_i$ there has
strictly positive valuation. Consequently $e_i^M h$, for a sufficiently
large fixed integer $M$, has positive coefficient valuation on all
those children. On the chosen child $e_i^M=1+$ positive-order terms,
so the existing reductions are unchanged. There are finitely many
nodes, and this constructs the required polynomial. A node with just
one leaf requires no filter. This also explains why independently
normalized Lagrange polynomials, with no control of their growth on
other scales, would not suffice.
\end{proof}

\begin{theorem}[Formal stationary balls]\label{thm:stationaryballs}
Let $\varphi(r,x)\in\C((r))[x]$ have degree $D\geq2$, and let
$h(r)\in\C((r))\setminus\{0\}$. Put
\[
 M=\C((r))[x]/(\partial_x+\varphi_x)\C((r))[x],\qquad
 \nabla[q]=[q_r+(\varphi_r+h'/h)q].
\]
After a finite ramification, this module has a full frame of rank-one
formal exponential functionals. If $a_D$ is the leading coefficient
of $\varphi$, every zero-exponential functional in this frame takes
the seed $[1]$ to a series of valuation at least
\begin{equation}\label{eq:valuationbound}
 v(h)-\frac{v(a_D)}D.
\end{equation}
In particular, if $Dv(h)>v(a_D)$, the differential submodule generated
by $[1]$ has no horizontal quotient taking $[1]$ to the constant one.
\end{theorem}
\begin{proof}
The two twisted derivatives commute, so the parameter connection
is well defined on the quotient. A nonzero polynomial of degree $m$
has spatial twisted derivative of degree $m+D-1$. Reduction of
successive leading terms therefore gives a basis
$1,x,\ldots,x^{D-2}$ and dimension $D-1$.

Split $\varphi_x$, retaining root multiplicities. For a critical root
$\xi$, expand
\[
 \varphi(\xi+u)=c_\xi+\sum_{j=2}^D a_{\xi,j}u^j,
 \qquad
 \rho_\xi=\max_{2\leq j\leq D,\ a_{\xi,j}\ne0}
                    \frac{-v(a_{\xi,j})}{j}.
\]
Enlarge the ramification so that all these scales are integral, and
set $\delta_\xi=r^{\rho_\xi}$. The polynomial
\[
 \varphi(\xi+\delta_\xi y)-c_\xi
\]
has integral coefficients, no linear term, and at least one unit
coefficient. Its reduction $F_\xi(y)$ has degree $d_\xi\geq2$.
Here and below we may replace $r$ by a root of $r$ and renormalize
the valuation; every displayed strict inequality is unchanged.
As $a_{\xi,D}=a_D$,
\begin{equation}\label{eq:balllower}
 \rho_\xi\geq-v(a_D)/D.
\end{equation}

\emph{The balls partition all critical roots.}
If another critical root $\zeta$ satisfies
$v(\zeta-\xi)\geq\rho_\xi$, translate the integral scaled polynomial
by $(\zeta-\xi)/\delta_\xi$. All its coefficients remain integral,
and the coefficient of degree $d_\xi$ remains a unit: higher-degree
terms have positive valuation. The linear term vanishes because
$\zeta$ is critical. It follows from the definition of the maximum
that $\rho_\zeta=\rho_\xi$. Thus their balls are identical.
If two of the balls intersect, one contains the other; the same
argument makes them identical. Otherwise they are disjoint.

The scaled derivative has reduction $F_\xi'$ of degree $d_\xi-1$.
It therefore has exactly $d_\xi-1$ integral roots, with multiplicity.
Here is a direct check of that count: factor it in a Puiseux field.
An integral root contributes a linear factor to the reduction;
for a root of negative valuation, factor out that root and the
remaining factor reduces to a constant. Multiplicativity of the
coefficient valuation leaves a nonzero constant scalar. Hence the
degree of the reduction counts exactly the integral roots.
The ball $B$ thus contains $d_B-1$ critical roots. Since
$\deg\varphi_x=D-1$,
\begin{equation}\label{eq:ballrankcount}
 \sum_B(d_B-1)=D-1.
\end{equation}

\emph{A functional for every dimension in a ball.}
The fixed polynomial $F_B$ has twisted quotient
\[
 H_B=\C[y]/(\partial_y+F_B')\C[y]
\]
with basis $1,y,\ldots,y^{d_B-2}$. Let $\ell^0_{B,j}$ be its dual
basis. Write
\[
 \varphi(\xi_B+\delta_B y)=c_B+F_B(y)+E_B(r,y),
 \qquad E_B\in r\C[[r]][y].
\]
For a polynomial amplitude $q$ define
\begin{equation}\label{eq:ballfunctional}
 \ell_{B,j}(q)=h\delta_B e^{c_B}
   \ell^0_{B,j}\!\left(e^{E_B(r,y)}q(\xi_B+\delta_B y)\right).
\end{equation}
Each coefficient uses finitely many polynomial terms. The exponential
in this formula is purely formal. The product rule conjugates the
spatial twisted derivative to $\partial_y+F_B'$, so the functional
annihilates the defining relations of $M$.

We spell out the parameter check because it retains the moving volume.
At fixed $y$, put $w=\xi_B'+\delta_B'y$. Differentiating
\eqref{eq:ballfunctional}, then subtracting its value on
$q_r+(\varphi_r+h'/h)q$, leaves, inside the fixed functional,
\[
 e^{E_B}\left(wq_x+\frac{\delta_B'}{\delta_B}q+\varphi_xwq\right)
 =\frac1{\delta_B}(\partial_y+F_B')
                      \left(e^{E_B}wq\right).
\]
It is zero. Thus the functionals are horizontal, including the
Jacobian $\delta_B$ and normalization $h$.

\emph{Completeness, not just a rank count.}
Choose the polynomials $q_{B,j}$ of
Lemma~\ref{lem:ballseparation}, for $0\leq j<d_B-1$.
Divide each evaluation column by its scalar $h\delta_Be^{c_B}$.
The resulting matrix reduces to the identity: $E_B$ has positive
order and the reductions of the amplitudes are the chosen monomials
on one ball and zero elsewhere. The matrix is therefore invertible.
Leading-degree reduction gives $\dim M=D-1$, so
\eqref{eq:ballrankcount} accounts for the entire module.

Remove from $c_B$ its finite negative-power part $q_B(r)$.
The remaining exponential is an invertible formal power series,
up to a nonzero constant. Formula~\eqref{eq:ballfunctional} and the
invertible matrix give the asserted direct sum of rank-one factors
$e^{q_B}$ times powers of $r$. If $q_B=0$, its value on the seed
has valuation at least $v(h)+\rho_B$, proving
\eqref{eq:valuationbound}.

Finally, after ramification this is a semisimple differential module.
A quotient on the submodule generated by the seed extends to the
whole module. Nonzero exponential factors have no horizontal map
to the Laurent field, because the logarithmic derivative of a
Laurent series cannot be the derivative of a nonzero negative-power
exponent. On the zero part, every horizontal functional is a
constant combination of those in the full frame. If the bound is
strictly positive, their seed values cannot sum to one.
\end{proof}

The construction is a formal stationary-phase argument. Its classical
background includes Newton--Puiseux and the formal stationary-phase
theory of exponential connections; Sabbah's local Fourier--Laplace
formula is a primary reference for the latter~\cite{SabbahStationary}.
We do not identify that formula's scope with the particular proof
above or claim priority for this formulation. The proof here also
allows nonconvergent formal Laurent coefficients. It uses a valued
field, not an arbitrary coefficient ring.

\begin{corollary}[Leading-coefficient obstruction]\label{cor:leadingobstruction}
For a rational polynomial phase $\varphi(t,x)$ of degree $D\geq2$
and nonzero rational scalar normalization $h(t)$, let $v_\infty$
mean order at $\varepsilon=1/t=0$. If
\begin{equation}\label{eq:leadingcriterion}
 v_\infty(h^D/a_D)>0,
\end{equation}
its seed submodule cannot have a rational horizontal quotient with
seed value one. In particular any horizontal formal functional at
another point cannot take that seed to one.
\end{corollary}
\begin{proof}
Apply the theorem at infinity. A rational horizontal quotient remains
a quotient on extension to the formal Laurent field. A formal
functional with seed value one restricts to such a rational quotient:
its values on successive seed derivatives are $1,0,0,\ldots$.
\end{proof}

The inequality must be strict. For example
$\varphi=-x^2/(2t^2)$ and $h=1/t$ have
$v_\infty(h^2/a_2)=0$, and their normalized Gaussian period is a
nonzero constant. Thus the earlier adverse scaling example is also
an exact boundary test of the criterion.


\subsection{A polynomial phase with logarithmic volume}\label{sec:puncturedpoly}

The puncture itself can be retained in the formal construction when
the phase is polynomial. This gives a stronger bound than the one for
$dx$: logarithmic volume does not change under a scaling of $x$.

\begin{theorem}\label{thm:puncturedpoly}
Let $K=\C((r))$, $\varphi\in K[x]$, and $h\in K\setminus\{0\}$.
On
\[
 H=K[x,x^{-1}]/(x\partial_x+x\varphi_x)K[x,x^{-1}],\qquad
 \nabla[q]=[q_r+(\varphi_r+h'/h)q],
\]
there is, after finite ramification, a full frame of rank-one formal
exponential functionals. Every zero-exponential functional in this
frame has seed value of valuation at least $v(h)$. Thus if $v(h)>0$,
the differential seed submodule has no horizontal quotient taking
$[1]$ to one. This includes a constant phase.
\end{theorem}
\begin{proof}
Write $D=\deg\varphi$, taking $D=0$ for a constant phase, including
zero. If $D\geq1$, the quotient has basis
$1,x,\ldots,x^{D-1}$. Indeed, on a negative monomial the lowest term
of $x\partial_x+x\varphi_x$ is its nonzero integer exponent times
that monomial. This eliminates all negative powers without introducing
lower ones. On polynomials, successive leading terms eliminate all
degrees at least $D$. Conversely an exact polynomial cannot come
from a Laurent polynomial with a negative lowest term, and a nonzero
polynomial preimage of degree $m$ has exact degree $m+D$.
This proves independence. For $D=0$, the quotient of $x\partial_x$
has the single basis element $1$.

Use the critical balls constructed in Theorem~\ref{thm:stationaryballs}.
At most one contains the origin. Every other ball has
$\rho_B>v(\xi_B)$, so $\gamma_B=\delta_B/\xi_B$ has positive
valuation. With the same $c_B,F_B,E_B$ and fixed functionals as before,
the local logarithmic functionals are
\begin{equation}\label{eq:puncturedaway}
 q\longmapsto h e^{c_B}\frac{\delta_B}{\xi_B}
  \ell^0_{B,j}\!\left(
      \frac{e^{E_B}q(\xi_B+\delta_B y)}{1+\gamma_B y}\right).
\end{equation}
The geometric denominator and all negative powers in $q$ expand
coefficientwise because $v(\gamma_B)>0$. Every coefficient is a
polynomial in $y$. The normalized logarithmic Jacobian is
$dx/x=\delta_B\,dy/(\xi_B+\delta_B y)$.

We retain the puncture in a separate central ball. If a critical
ball contains zero, write it as $v(x)\geq\rho_0$ and put
$\delta_0=r^{\rho_0}$. Translation within that ball gives
\[
 \varphi(\delta_0 y)=c_0+F_0(y)+E_0(r,y),\qquad
 c_0=\varphi(0),\quad E_0\in r\C[[r]][y],\quad F_0(0)=0.
\]
Here $d_0=\deg F_0\geq2$, and the ball contains $d_0-1$
critical roots with multiplicity. If no critical ball contains zero,
choose $\rho_0$ sufficiently large that the central ball is disjoint
from all critical balls and every nonconstant coefficient of
$\varphi(\delta_0 y)$ has positive valuation. Then $F_0=0$.
Such a choice exists: there are finitely many roots and coefficients.
This also handles $D\leq1$ without any critical balls.

For $d_0\geq2$, the fixed punctured quotient
\[
 H_0=\C[y,y^{-1}]/(y\partial_y+yF_0')\C[y,y^{-1}]
\]
has basis $1,y,\ldots,y^{d_0-1}$, by the same reduction argument.
For $F_0=0$, it has basis $1$ and dual functional the constant term.
For each element $\lambda_j$ of the dual basis define
\begin{equation}\label{eq:puncturedcentral}
 q\longmapsto h e^{c_0}\lambda_j
                   \!\left(e^{E_0}q(\delta_0 y)\right).
\end{equation}
There is no factor $\delta_0$: $dx/x=dy/y$. Each coefficient is
a Laurent polynomial, so the expression is well defined.

We check both identities, including parameter dependence. Spatial
exactness in the central ball is conjugation by $e^{E_0}$ of
$y\partial_y+yF_0'$. The derivative discrepancy is
$\delta_0'/\delta_0$ times the same exact expression applied to
$e^{E_0}q$. On an outside ball put $x=\xi_B+\delta_B y$ and
$w=\xi_B'+\delta_B'y$. After removing $h e^{c_B}$, the derivative
discrepancy in \eqref{eq:puncturedaway} is
\[
 e^{E_B}\left[
   \left(\frac{\delta_B'}x-\frac{\delta_B w}{x^2}\right)q
        +\frac{\delta_B w}x(q_x+\varphi_xq)\right]
 = (\partial_y+F_B')\left(e^{E_B}\frac wxq\right).
\]
It is annihilated by $\ell^0_{B,j}$. The identical product rule
proves spatial exactness there. Thus every displayed functional is
horizontal for the stated connection.

Completeness follows from polynomial separation of these disjoint
balls, not merely their number. For the central ball use the monomials
of the basis of $H_0$; for each outside ball use
$1,\ldots,y^{d_B-2}$. Lemma~\ref{lem:ballseparation} gives polynomial
amplitudes with exactly these reductions on their own ball and zero
on the others. Normalize the central columns by $h e^{c_0}$ and
the outside columns by $h e^{c_B}\delta_B/\xi_B$. Their matrix
reduces to the identity. If the central ball contains critical roots,
it contributes their multiplicity plus one; otherwise it contributes
one. Together with all outside roots this gives $D$ dimensions for
$D\geq1$, or one for $D=0$, the full quotient rank.

As in Theorem~\ref{thm:stationaryballs}, this frame is a direct sum
of rank-one exponential factors after ramification. For a zero factor,
\eqref{eq:puncturedcentral} has seed valuation at least $v(h)$;
\eqref{eq:puncturedaway} has valuation at least
$v(h)+\rho_B-v(\xi_B)>v(h)$. Semisimplicity now extends a quotient
of the differential seed submodule to the full module. Nonzero
exponential factors cannot map horizontally to the Laurent field.
A constant combination of the zero-factor seed values cannot be
one if $v(h)>0$.
\end{proof}

\begin{remark}[Why two irregular ends require another argument]
For $\varphi=r(x+x^{-1})$ and $h=1$, put $\theta=r\partial_r$.
Every logarithmic period satisfies
\[
 (\theta^2-4r^2)I=0.
\]
Indeed $y=x+x^{-1}$ and $z=x-x^{-1}$ obey
$(x\partial_x+x\varphi_x)z=y+rz^2$ and $y^2-z^2=4$.
These identities give the equation by integration by parts.
The quotient has rank two, and $[1],[y]$ are independent:
its basis is $x^{-1},1$ and $[y]=2[x^{-1}]$.

An explicit formal basis of scalar solutions is
\[
 I=\sum_{n\geq0}\frac{r^{2n}}{(n!)^2},\qquad
 J=I\log r-\sum_{n\geq1}\frac{H_n r^{2n}}{(n!)^2},
 \quad H_n=\sum_{j=1}^n\frac1j.
\]
The coefficient recurrence $n^2a_n=a_{n-1}$ proves the first
equation; $H_n-H_{n-1}=1/n$ proves the second.
Their $\theta$-Wronskian is one: its derivative is zero by the
equation and its constant term is one. Thus the logarithm is a full
rank-two phenomenon, with nontrivial unipotent monodromy that no
finite ramification removes. These are the familiar modified Bessel
solutions; see~\cite[10.25.1, 10.31.2]{DLMF}.
This example does not refute the desired logarithmic bound:
multiplication by $h=r$ makes both seed values tend to zero.
It does refute an extension of the semisimple polynomial-frame theorem
to all Laurent phases. The remaining proof must allow such extensions.
\end{remark}

The residue dimension at the puncture is essential. Dropping it would
give rank $D-1$ instead of $D$ and an incomplete frame. This theorem
still permits nonconvergent formal coefficients, but is a valued-field
statement. It does not cover a phase with both positive and negative
powers of $x$.

\begin{corollary}\label{cor:oneend}
Suppose $\varphi(t,x)$ is rational in $t$ and polynomial either in
$x$ or in $x^{-1}$. For the logarithmic-volume connection normalized
by $h=1/t$, no horizontal formal functional at another point can have
seed value one.
\end{corollary}
\begin{proof}
For a polynomial in $x$, apply Theorem~\ref{thm:puncturedpoly} at
$r=1/t=0$, where $v(h)=1$. The rational seed restriction is justified
as in Corollary~\ref{cor:leadingobstruction}. For a polynomial in
$x^{-1}$ use $y=1/x$; $dx/x=-dy/y$ only changes a constant factor.
\end{proof}

\section{Logarithmic Gaussian charts}
Let $h(Z,T)=a+bZ+cT+dZ^2+eZT+fT^2$ and
$\Delta=e^2-4df\ne0$, and let $P=Wh(Z,T)+G(Z,T)$ with
arbitrary polynomial $G$. The coordinate charts below and their
formal transport retain the actual small circular fixed point.
\subsection{The logarithmic branch, including its removed component}

For $\Delta\ne0$, $f\ne0$, choose $\eta^2=\Delta$ and set
\[
 \rho=\frac{e+\eta}{2f},\quad \sigma=\frac{e-\eta}{2f},\quad
 u=\frac{\rho c-b+1/t}{\eta},\quad
 v=\frac{b-1/t-\sigma c}{\eta},\quad k=uv-a/f.
\]
On the circular constraint $z=th(z,T)$, put
$x=T+\rho z+u$, so that $T+\sigma z+v=k/x$. Then
\[
 z=\frac f\eta(x-k/x-u+v),\quad T=x-u-\rho z,
 \qquad
 \frac{dT}{1-th_Z}=\frac{dx}{t\eta x}.
\]
The coefficient of $t^{-2}$ in $k$ is $-1/\eta^2$, so the conic
is generically smooth. The exponent $tG(z,T)-T^2/2$ is Laurent
polynomial in $x$. On the formal Gaussian branch,
$x=1/(\eta t)+O(1)$, which makes every Laurent power legitimate.

If $f=0$, then $e\ne0$. Put $z_0=-c/e$ and $x=z-z_0$. With
$A(z)=a+bz+dz^2$, the constraint gives
\[
 T=\frac{z/t-A(z)}{ex},\qquad
 \frac{dT}{1-th_Z}=\frac{dx}{tex}.
\]
This localization is valid if $c\ne0$, or if $c=0$ and $a\ne0$.
If $a=c=0$, the chosen fixed point is instead $g=0$ identically;
inverting $g$ would discard the very branch being tested. In that
case the exact unit series is
\[
 \E_T\left[\frac{e^{tG(0,T)}}{1-t(b+eT)}\right].
\]
Use $x=T-(1/t-b)/e$, now a valid Laurent unit. Its volume is
$-dx/(tex)$ and its exponent is polynomial in $x$. These formulas
cover both the smooth conic and the removed-component branch.
We next connect these coordinates to the formal Gaussian functional.
This will close every one-ended phase without an analytic estimate.

\begin{lemma}[Formal logarithmic transport]\label{lem:logtransport}
In each of the three charts above, write $X=X(t,T)$ for the formal
coordinate, $D$ for the displayed circular denominator, and $B$ for
$tG(g,T)$ or $tG(0,T)$ on the removed component. Put
$\kappa=\eta,e,-e$, respectively. Then
\[
 J(q)=\E_T\!\left[\frac{q(X)e^B}{D}\right]
\]
is horizontal on the logarithmic twisted quotient with phase
$\varphi$ and normalization $h=1/t$, and $J(1)=\Phi_1$.
\end{lemma}
\begin{proof}
All Laurent powers of $X$ are legitimate by the unit checks above.
The coordinate identities give
\[
 X_T=t\kappa X/D,\qquad
 \varphi(t,X)=B-T^2/2.
\]
Consequently Gaussian integration by parts gives
$J((x\partial_x+x\varphi_x)q)=0$: the corresponding integrand,
multiplied by $t\kappa$, is
$(\partial_T-T)\{q(X)e^B\}$.

For parameter horizontality set $w=1/D$ and $R=X_t/X_T$.
Differentiating $w=X_T/(t\kappa X)$ gives the exact identity
\[
 w_t+w/t=\partial_T(wR).
\]
Subtracting $J(q_t+(\varphi_t-1/t)q)$ from $\partial_tJ(q)$ now
gives the Gaussian expectation of
$(\partial_T-T)\{wR q(X)e^B\}$, hence zero. These are formal
coefficientwise identities, without an integral-convergence hypothesis.
The seed equality is the circular contraction formula.
\end{proof}

\begin{proposition}\label{prop:oneendgaussian}
Every chart above whose actual Laurent phase has at most one
irregular end is not pure. The removed-component branch
$a=c=f=0$, $e\ne0$, is therefore excluded for every polynomial $G$,
without a degree bound.
\end{proposition}
\begin{proof}
At most one irregular end means that $\varphi$ is polynomial in
$x$ or in $x^{-1}$, including the constant case. Purity would give
the horizontal functional in Lemma~\ref{lem:logtransport} seed value
one, contrary to Corollary~\ref{cor:oneend}. The removed-component
phase is always polynomial in $x$, irrespective of $\deg G$.
\end{proof}

Thus only phases with genuinely nonzero positive and negative Laurent
orders remain in the logarithmic estimate. Their coefficient
specializations must be judged using the actual orders, not generic
ones. The formal proof above does not establish a spanning flat cut
family for two irregular ends.


\begin{thebibliography}{9}
\bibitem{SabbahStationary}
C.~Sabbah, \emph{An explicit stationary phase formula for the local
formal Fourier--Laplace transform}, in \emph{Singularities I},
Contemporary Mathematics \textbf{474}, American Mathematical Society,
2008, 309--330;
\href{https://arxiv.org/abs/0706.3570v2}{arXiv:0706.3570v2}.

\bibitem{Wall}
C.~T.~C.~Wall, \emph{Singular Points of Plane Curves}, Cambridge
University Press, 2004, Chapter~2, 15--38,
\href{https://doi.org/10.1017/CBO9780511617560.003}{doi:10.1017/CBO9780511617560.003}.

\bibitem{DLMF}
NIST Digital Library of Mathematical Functions,
\href{https://dlmf.nist.gov/10.25.E1}{10.25.1} and
\href{https://dlmf.nist.gov/10.31.E2}{10.31.2},
modified Bessel equation and logarithmic power series.

\end{thebibliography}
\end{document}
